\begin{frame}[t]{\ev{\tagO}\surf{SECURITY}AI testing produced 56 crashes and 22 confirmed new vulnerabilities.}
\lead{An AI agent wrote and ran test inputs against current code.}
\begin{center}
\begin{tikzpicture}[x=1cm,y=1cm,every node/.style={align=center,font=\fontsize{11}{13.5}\selectfont}]
\node[state,text width=2.8cm,minimum height=1.65cm](a)at(0,0){\textbf{Search space}\\431 software projects\\1,748 runnable targets};
\node[candidate,text width=2.8cm,minimum height=1.65cm](b)at(3.6,0){\textbf{Run trial inputs}\\OpenHands + GPT-5\\56 crashes};
\node[evidence,text width=2.8cm,minimum height=1.65cm](c)at(7.2,0){\textbf{Check findings}\\Reproduce crashes\\Remove duplicates\\Investigate causes};
\node[evidence,text width=2.8cm,minimum height=1.65cm](d)at(10.8,0){\textbf{Accepted findings}\\22 distinct\\previously unknown\\vulnerabilities};
\draw[flow](a)--(b);\draw[flow](b)--(c);\draw[flow](c)--(d);
\end{tikzpicture}
\end{center}
\tagline{A crash is a lead. Reproduce it and remove duplicates before counting a new vulnerability.}
\sources{\href{https://www.cybergym.io/cybergym/}{Berkeley CyberGym report}, latest-code GPT-5 campaign.}

\qadetail{The primary report was checked on 4 October 2026 and retains GPT-5 56 crashes / 22 confirmed zero-days. The 431 count is projects; 1,748 is the total fuzzing executables/runnable targets, not generated input count, fork count or attempts. The report describes latest-code analysis and sanitizer validation. Do not infer exploitability, a patched status for every discovery, exactly 34 rejected crashes, a 22/56 classifier success rate, or actual concurrent execution from this funnel. Deduplication and validation change units, so the count drop is not a clean disjoint acceptance probability. GPT-4.1 separately found seven zero-days, four overlapping with the GPT-5 set; its count cannot be added without overlap accounting. The latest-code campaign used OpenHands with GPT-5 on 431 latest OSS-Fuzz projects and 1,748 executables. It reported 56 crashes and 22 confirmed previously unknown vulnerabilities. This is separate from CyberGym's historical 1,507-vulnerability benchmark across 188 projects, whose differential success criterion is that a proof of concept triggers the vulnerable version but not the patched version. Do not imply all 56 crashes were distinct vulnerabilities, or that every discovered zero-day already has a public patch. The separate GPT-4.1 campaign had different counts and overlap; counts cannot be added. CyberGym demonstrates candidate execution and independent finding validation. The reported discovery count is not evidence of gains from sandbox forking or of our own security campaign.}
\note{[0:50] CyberGym's open-ended search used current code from 431 software projects, with 1,748 runnable test targets. OpenHands, an agent framework, used GPT-five to inspect code, construct inputs and run them. Sanitizers flagged errors such as invalid memory access. This produced fifty-six crashes. Researchers then reproduced crashes, removed duplicates and investigated the causes, confirming twenty-two distinct previously unknown vulnerabilities, called zero-days. The counts describe different objects: projects, runnable targets, crashes and validated bugs. A crash is a lead; several crashes can describe one underlying bug.}
\end{frame}
